% ======================================================================
% CAPITULO 15 — Ciclos Evolutivos do Raciocinio Cientifico
% ======================================================================
\chapter{Ciclos Evolutivos do Raciocinio Cientifico}

\section{O Que E um Ciclo Evolutivo de Raciocinio?}

No OpenCode Ecosystem, o desenvolvimento nao e linear --- e \textbf{ciclico}.
Cada experimento, cada artigo, cada validacao gera aprendizado que e
incorporado ao sistema, tornando-o progressivamente mais capaz. Este capitulo
documenta como o raciocinio cientifico evolui atraves de ciclos, usando GAT
como caso de estudo\footnote{Este capitulo e assumidamente metacientifico:
estamos usando o OpenCode para estudar como o proprio OpenCode evolui seu
raciocinio. Conforme o Principio de Integridade (R-I5: Contraprova
Independente), distinguimos claramente entre o que e auto-observacao do
sistema e o que e validacao externa.}.

\subsection{A Metafora do Cientista em Treinamento}

\begin{center}
\fbox{\begin{minipage}{0.88\textwidth}
\textbf{Para leigos:} Pense num estudante de doutorado. No primeiro ano, ele
mal consegue resolver os exercicios do livro. No segundo, ja escreve um artigo
aceitavel. No quarto, defende uma tese original. O que mudou? Nao foi o QI
--- foi o \textbf{processo}: a cada ciclo (ler $\to$ entender $\to$ aplicar
$\to$ errar $\to$ corrigir $\to$ publicar), o raciocinio se refina. O OpenCode
automatiza exatamente esse ciclo, acelerando-o de anos para minutos.
\end{minipage}}
\end{center}

\section{Os 17 Ciclos do OpenCode (2026)}

O ecossistema passou por 17 ciclos de desenvolvimento documentados.
Cada ciclo e uma iteracao completa de \textbf{PLAN $\to$ ACT $\to$ REFLECT
$\to$ EXTRACT $\to$ EVOLVE}\footnote{O ciclo PlanAct e inspirado no
framework ReACT de agentes de IA: Yao et al. (2023). \textit{ReAct: Synergizing
Reasoning and Acting in Language Models}. ICLR. DOI:
10.48550/arxiv.2210.03629. \textbf{Para leigos:} Em vez de ``pensar e depois
agir'', o agente alterna entre pensar e agir --- como um aluno que resolve um
problema passo a passo, verificando cada etapa antes de prosseguir.}.
Abaixo, os 5 ciclos mais relevantes para a GAT:

\subsection{Ciclo 1 --- Cross-Validation Inicial (Score: 85/100)}

\textbf{O que foi feito:} Primeira integracao com dados reais (World Bank API).
27 indicadores de 50 paises analisados com bootstrap de 10.000 reamostragens.

\textbf{Descoberta relevante para GAT:} A correlacao entre educacao e inovacao
foi $r = -0,03$ (essencialmente zero), enquanto P\&D privado mostrou $r = 0,73$.
Isso ilustra um principio fundamental: \textbf{nem toda variavel observavel e
um fator de risco precificavel} --- assim como nem toda flutuacao de preco
constitui arbitragem.

\textbf{Licao para o modulo:} A distincao entre correlacao e causalidade e
central tanto para a econometria quanto para a GAT. Um mercado pode ter
curvatura zero (sem arbitragem) e ainda assim apresentar correlacoes entre
ativos --- a estrutura de correlacao e capturada pela metrica do fibrado, nao
pela curvatura.

\subsection{Ciclo 6 --- Correcao Iterativa (Score: 86,5 $\to$ 92,7)}

\textbf{O que foi feito:} Implementacao de banca simulada com 5 revisores e
4 orientadores doutores. Loop de correcao: submeter $\to$ revisar $\to$
corrigir $\to$ reavaliar $\to$ repetir ate score $\geq$ 95.

\textbf{Licao para GAT:} O processo de correcao iterativa e \textbf{identico
em estrutura} ao processo de eliminacao de arbitragem em mercados reais:

\[
\begin{array}{c|c}
\textbf{Correcao Academica} & \textbf{Mercado Financeiro} \\ \hline
\text{Submeter artigo} & \text{Precificar ativo} \\
\text{Revisor encontra erro} & \text{Arbitrador encontra inconsistencia} \\
\text{Autor corrige} & \text{Preco se ajusta} \\
\text{Revisor reavalia} & \text{Oportunidade desaparece} \\
\text{Artigo aprovado (score $\geq$ 95)} & \text{Mercado em equilibrio (R = 0)}
\end{array}
\]

\subsection{Ciclo 9 --- SDD+TDD (Score: 94/100)}

\textbf{O que foi feito:} 7 especificacoes modulares, 9 criterios de teste,
7 correcoes aplicadas, 3 ADRs DecisionNode, 16 perguntas de banca simuladas.

\textbf{Licao para GAT:} O ciclo SDD+TDD e o analogo metodologico do
\textbf{Principio da Curvatura} da GAT:

\begin{enumerate}[nosep]
  \item \textbf{SPEC} (especificacao) = definir o espaco de carteiras $M$
  \item \textbf{CT} (criterios de teste) = definir a conexao $\chi$
  \item \textbf{TDD} (testes) = verificar se a curvatura $R$ e zero
  \item \textbf{AUDITORIA} = verificar se Hol$^0(\chi)$ e trivial
\end{enumerate}

Se o TDD falha ($R \neq 0$), ha um ``bug'' no modelo --- exatamente como
ha arbitragem no mercado. Corrigir o modelo (REFACTOR) e analogo a eliminar
a oportunidade de arbitragem.

\subsection{Ciclo 16 --- Teoria dos Jogos (Score: 96/100)}

\textbf{O que foi feito:} Integracao de 10 estrategias de teoria dos jogos ao
orquestrador de raciocinio: Nash, Minimax, Shapley, Harsanyi, Bayesian-Nash,
Evolutivo, Coalizao, Barganha, Sinalizacao, Leilao.

\textbf{Licao para GAT:} A conexao entre Teoria dos Jogos e GAT e profunda,
mas nao obvia\footnote{Nash, J.F. \textit{Equilibrium points in n-person games}.
PNAS, 36(1):48--49, 1950. DOI: 10.1073/pnas.36.1.48. \textbf{Resenha para
leigos:} John Nash (o genio retratado no filme ``Uma Mente Brilhante'')
provou que em qualquer jogo com um numero finito de jogadores e estrategias,
existe pelo menos um ponto de equilibrio --- uma situacao onde ninguem se
arrepende da estrategia escolhida. Este conceito rendeu-lhe o Nobel de
Economia em 1994.}:

\begin{itemize}[nosep]
  \item \textbf{Equilibrio de Nash:} Um mercado em equilibrio (NFLVR, $R=0$)
        e um equilibrio de Nash do jogo entre arbitradores e formadores de
        mercado.
  \item \textbf{Valor de Shapley:} A contribuicao marginal de cada ativo para
        o risco sistemico da carteira --- analogo a decomposicao da curvatura
        por direcao no fibrado.
  \item \textbf{Jogos Bayesianos (Harsanyi):} Mercados com informacao
        incompleta --- agentes nao conhecem todos os precos. A curvatura pode
        ser nao-nula simplesmente porque a informacao e assimetrica.
\end{itemize}

\subsection{Ciclo 17 --- CORA-Eval Benchmark (Score: 97/100)}

\textbf{O que foi feito:} 150 tarefas em 10 dimensoes $\times$ 4 niveis
(Basico a Pesquisa). CORA-Score 3.04 (M4). Validacao externa: 34/34
Project Euler + Rosalind.

\textbf{Licao para GAT:} O CORA-Eval e o analogo, para raciocinio cientifico,
do que o \textbf{Teorema Central da GAT} e para arbitragem: uma metrica que
captura a ``maturidade'' do raciocinio. Assim como a curvatura $R$ mede o
grau de arbitragem, o CORA-Score mede o grau de sofisticacao cientifica.

\section{A Evolucao do Raciocinio sobre GAT}

Aplicando o metodo SDD+TDD ao estudo da propria GAT, documentamos como o
raciocinio evoluiu ao longo dos ciclos:

\begin{longtable}{p{0.05\textwidth}p{0.25\textwidth}p{0.25\textwidth}p{0.25\textwidth}}
\toprule
\textbf{C} & \textbf{Pergunta Inicial} & \textbf{Hipotese} & \textbf{Resposta apos Verificacao} \\
\midrule
1 & Existe uma linguagem geometrica para arbitragem? & Sim, via fibrados principais & Confirmado: $R \neq 0 \Leftrightarrow$ arbitragem. Mas $R=0 \not\Rightarrow$ NFLVR (Novikov) \\
2 & A holonomia classifica estrategias? & Sim, cada classe $\leftrightarrow$ elemento de Lie(Hol) & Parcialmente confirmado. A classificacao e valida para mercados sem friccoes. \\
3 & O mercado tem um Lagrangiano? & Sim, $\mathcal{L}(D, \mathcal{D}D, t)$ & Confirmado. Solucoes de equilibrio correspondem a pontos criticos. \\
4 & A arbitragem satisfaz uma equacao de continuidade? & Sim, div $J = 0 \Leftrightarrow$ NFLVR & Confirmado. A analogia hidrodinamica e matematicamente precisa. \\
5 & E possivel implementar GAT computacionalmente? & Sim, via Python + numpy & Confirmado. Curvatura calculavel para mercados com $N \leq 100$ ativos. \\
\bottomrule
\caption{Evolucao do raciocinio sobre GAT ao longo dos ciclos}
\label{tab:evolucao_gat}
\end{longtable}

\section{Protocolo de Experimento Cientifico com OpenCode}

Para cada experimento GAT no OpenCode, siga este protocolo:

\begin{enumerate}[nosep]
  \item \textbf{SPEC:} Escreva a especificacao do experimento (hipotese,
        metodo, criterios de sucesso). Salve em \texttt{specs/}.
  \item \textbf{TDD:} Escreva o teste automatizado que verifica a hipotese.
        O teste deve falhar inicialmente (RED).
  \item \textbf{IMPLEMENTAR:} Escreva o codigo do experimento.
  \item \textbf{EXECUTAR:} Rode o teste. Corrija ate passar (GREEN).
  \item \textbf{AUDITAR:} Submeta ao Cora-Debate V1--V7.
  \item \textbf{DOCUMENTAR:} Gere o relatorio com \texttt{opencode run /artigo}.
  \item \textbf{EVOLUIR:} O AutoEvolve analisa o ciclo e gera novas skills.
  \item \textbf{REPETIR:} Volte ao passo 1 com uma pergunta mais avancada.
\end{enumerate}

\subsection{Exemplo Completo: Verificar o Teorema Central}

\begin{lstlisting}[caption={Protocolo completo: verificar Teorema Central}]
# SPEC-GAT-002: Verificar que R=0 ⇔ existe pricing kernel

# PASSO 2: TDD (teste antes do codigo)
def test_curvatura_zero_implica_pricing_kernel():
    """CT-1: Se R=0, entao existe beta>0 com d(beta*D) martingal"""
    mercado = criar_mercado_ito(N=3, mu=[0.03, 0.03, 0.03], r=0.03)
    R = calcular_curvatura(mercado)
    assert np.allclose(R, 0), f"R deveria ser zero, obtido {R}"
    beta = encontrar_pricing_kernel(mercado)
    assert beta is not None, "Pricing kernel deveria existir"
    assert all(beta > 0), "Pricing kernel deve ser positivo"

# PASSO 3-4: Implementar e verificar
# (codigo em evaluations/tests/test_gat_teorema_central.py)

# PASSO 5: Auditar com Cora-Debate
# opencode run /auditar --spec SPEC-GAT-002

# PASSO 6-7: Documentar e evoluir
# opencode run /artigo --tema "Verificacao do Teorema Central da GAT"
\end{lstlisting}

\section{Exercicios}

\begin{exer}[Seu Primeiro Ciclo Evolutivo]
Escolha um topico da GAT que voce achou confuso neste modulo. Siga o
protocolo de 8 passos para transforma-lo em um experimento verificado.
Documente:
\begin{enumerate}[(a)]
  \item Qual era sua duvida inicial?
  \item Que hipotese voce formulou?
  \item O que o experimento revelou?
  \item O que voce aprendeu que nao sabia antes?
\end{enumerate}
\end{exer}

\begin{exer}[Comparacao entre Ciclos]
Compare os resultados do Ciclo 1 (cross-validation inicial) com o Ciclo 17
(CORA-Eval). O que mudou na \textbf{qualidade} do raciocinio? Que metricas
capturam essa mudanca? Como voce aplicaria esse conceito de ``maturidade de
raciocinio'' ao estudo da GAT?
\end{exer}

\begin{exer}[Projete o Ciclo 18]
Com base no que voce aprendeu sobre GAT e sobre o OpenCode, proponha o
Ciclo 18. Que pergunta ele deveria responder? Que experimento deveria
executar? Que nova skill deveria gerar? Justifique com base na Secao 15.3.
\end{exer}

\section{Referencias do Capitulo}

\begin{enumerate}[nosep]
  \item Claro, M. (2026). OpenCode Ecosystem v4.6.1. GitHub.
  \item Yao, S. et al. (2023). ReAct: Synergizing Reasoning and Acting in
        Language Models. ICLR. DOI: 10.48550/arxiv.2210.03629.
  \item Nash, J.F. (1950). Equilibrium points in n-person games. PNAS.
        DOI: 10.1073/pnas.36.1.48.
  \item Beck, K. (2003). Test-Driven Development: By Example. Addison-Wesley.
  \item Popper, K. (1959). The Logic of Scientific Discovery. Routledge.
        DOI: 10.4324/9780203994627.
\end{enumerate}
